%! Characteristics of Logarithmic Functions — Unit 8, Lesson 8.0 — Student Packet Cover Sheet
\documentclass[10pt]{article}
\usepackage{algebra2-article}
\usepackage{algebra2-boxes}
\usepackage{ltablex}
\keepXColumns

\begin{document}

% Full-bleed forest banner
\begin{tikzpicture}[remember picture, overlay]
  \fill[forest] (current page.north west)
    rectangle ([yshift=-0.9in]current page.north east);
\end{tikzpicture}
\vspace{-0.8in}

\begin{center}
  {\color{white}\textbf{\LARGE Algebra 2: Shepherd}}\\[4pt]
  {\color{forestmid}\large\textbf{Unit 8 \quad Logarithmic Functions}}\\[6pt]
  {\color{white}\textbf{\large Lesson 8.0 \quad Characteristics of Logarithmic Functions}}
\end{center}

\vspace{0.22in}
\namedateperiod
\vspace{0.18in}

\begin{learningtargetbox}
\small
\begin{enumerate}[leftmargin=1.5em, itemsep=5pt]
  \item \ldots{} find a logarithmic function's \textbf{domain} by setting its \textbf{argument} $>0$, and name the \textbf{vertical asymptote} that sits at the edge of that domain.
  \item \ldots{} read the \textbf{domain}, \textbf{range}, \textbf{asymptote}, \textbf{intercepts}, \textbf{increasing/decreasing} and \textbf{positive/negative} intervals, \textbf{extrema}, and \textbf{end behavior} off a logarithmic graph with base $b>1$ or $0<b<1$.
  \item \ldots{} explain the \textbf{inverse relationship} between $y=\log_b x$ and $y=b^x$ --- the domain/range swap and the asymptote that turns a corner --- and \textbf{compare and contrast} the two families.
\end{enumerate}
\end{learningtargetbox}

\vspace{0.14in}

\begin{tocbox}
\small
\renewcommand{\arraystretch}{1.5}
\begin{tabularx}{\linewidth}{@{} c l X r @{}}
\rowcolor{forestbg}
\textbf{\#} & \textbf{Component} & \textbf{Description} & \textbf{Score} \\
\midrule
1 & Warm-Up        & Spiral review: exponents run backwards, the exponential parent, and a pair of traded tables & \blank{1.2cm} \\
2 & Guided Notes   & Reading every characteristic off a logarithmic graph, and where the graph came from & \blank{1.2cm} \\
3 & Group Activity & \emph{Reading Logarithmic Graphs} --- feature tables, domains from equations, and a doubling model, by tier & \blank{1.2cm} \\
4 & Exit Ticket    & Quick check: features from a graph, plus an SOL-style domain item & \blank{1.2cm} \\
5 & Homework       & Domains from equations, two contrasting graphs, and building a log from an exponential table & \blank{1.2cm} \\
\midrule
  & \multicolumn{2}{r}{\textbf{Total}} & \blank{1.2cm} \\
\end{tabularx}
\end{tocbox}

\vspace{0.10in}

\begin{remindbox}
\small A \textbf{logarithmic function} $y=\log_b x$ answers an exponential question backwards:
$\log_b x$ is \emph{the exponent you put on $b$ to get $x$}. Because no power of a positive base is
ever zero or negative, the \textbf{argument} of a logarithm must be \emph{strictly} positive --- so
to find a domain you set the argument $>0$ and solve. The input the domain stops at is not part of
the graph; it is a \textbf{vertical asymptote}, a line the curve approaches forever and never
reaches. This is the difference from Unit~6: a square root got to \emph{stand on} its endpoint, but a
logarithm's edge is missing. Every parent log graph has range \emph{all reals}, an $x$-intercept at
$(1,0)$ (since $b^0=1$ for any base), \emph{no} $y$-intercept, and \emph{no} extrema; a base $b>1$
makes it increase and a base $0<b<1$ makes it decrease. All of it comes from one fact: $y=\log_b x$
is the \textbf{inverse} of $y=b^x$, reflected over the line $y=x$. Inputs and outputs trade places,
so the domains and ranges swap --- and the exponential's \emph{horizontal} asymptote stands up into
the logarithm's \emph{vertical} one.
\end{remindbox}

\end{document}
